\begin{recipe}[
    preparationtime = {30min},
    bakingtime = {45min},
    bakingtemperature = {\protect\bakingtemperature{
        topbottomheat=\unit[180]{\textcelcius}
    }},
    portion = {\portion{8}},
     calory = {860kcal | 11g Protein | 55g KH | 67g Fett},
    price = {1,08€},
    rating = {3},
    source = {\protect\recipesource{https://kitchen-by-the-sea.com/tiramisu-cake-2/?fbclid=PAT01DUAN0r1hleHRuA2FlbQIxMAABpx8rGBCjVLX8xsF16bC1j2Lxq\_6TTiUbXqFM4BLas7-R9rcNzZxMg7GtTsZ6\_aem\_UsdUqGjVamN-gTxW6OO4Iw\#recipe}{Kitchen by the sea: Tiramisu Cake}}
]{Tiramisu Chocolate Cake}

\graph{% pictures
        big=pic/dessert/07_tiramisu-cake  % big picture
    }

\ingredients{
    \textbf{Kuchen} & \\
    90g & Mehl\\
    60g & Kakaopulver\\
    200g & Zucker\\
    1 TL & Natron\\
    \nicefrac{1}{2} TL & Backpulver\\
    \nicefrac{1}{2} TL & Salz\\
    60ml & Öl\\
    120ml & Buttermilch\\
    1 & Ei\\
    1 TL & Vanille\\
    120ml & Warmes Wasser oder Kaffee\\
    \textbf{Creme} & \\
    500g & Mascarpone\\
    480ml & Sahne\\
    65g & Zucker\\
    \nicefrac{1}{4} TL & Salz\\
    1 EL & Vanille\\
    2 EL & Instant Espresso\\
    3 EL & Frisch gemahlener Kaffee\\
    \textbf{Topping} & \\
    3 EL & Kakaopulver
}

\preparation{
    \step%
    Backofen auf 175°C vorheizen. 23cm Springform mit Backpapier auslegen.
    
    \step%
    Mehl, Kakao, Zucker, Natron, Backpulver und Salz vermengen.
    
    \step%
    Öl, Buttermilch, Ei, Vanille und Wasser zugeben und glatt rühren.
    
    \step%
    Teig in die Form füllen und 45min backen. Vollständig auskühlen lassen.
    
    \step%
    Für die Creme Mascarpone und Sahne 30s verrühren.
    
    \step%
    Zucker, Salz, Vanille und Espresso zugeben und bis zu soft peaks schlagen.
    
    \step%
    Gemahlenen Kaffee vorsichtig unterheben.
    
    \step%
    Creme auf den ausgekühlten Kuchen streichen und mindestens 4h (besser über Nacht) kühlen.
    
    \step%
    Mit Kakaopulver bestäuben.
}

\hint{Creme nicht übermixen, Mascarpone trennt sich sonst und wird körnig.}

\end{recipe}